\documentclass[10pt,a4paper]{article}
\usepackage[utf8]{inputenc}
\usepackage[margin=25mm]{geometry}

\usepackage[english]{babel}
\usepackage[style=archaeologie]{biblatex}


\usepackage{csquotes}
\usepackage{fontawesome} % cool icons. Usage: \faBook; the whole list of available ones:  ftp://ftp.dante.de/tex-archive/fonts/fontawesome/doc/fontawesome.pdf
\usepackage[default]{raleway} % chooses the raleway font instead of LaTeX's standard
\usepackage{hyperref}
\usepackage{graphicx}
\usepackage{epigraph} % for the dramatic moment or two
\usepackage{wrapfig} % for pictures beside the text

\setlength{\parindent}{0pt} % stops paragraph indentation (remove if you want)
\newcommand{\dotdot}{ \lbrack\dots\rbrack~} % custom command to get [...] for quoting
\setlength{\emergencystretch}{.25em} % avoid line overflow 
\clubpenalty10000\widowpenalty10000\displaywidowpenalty=10000% avoid widow lines and orphans

\bibliography{literature}
%--------------------------------
\title{A Humanities' seminar paper with \LaTeX{} -- in 10 minutes}
\author{\LaTeX{} Ninja}
\date{\today}
\begin{document}
\maketitle
\tableofcontents
\epigraph{All human things are subject to decay, and when fate  summons, Monarchs must obey}{\textit{Mac Flecknoe \\ % \\ causes manual line break
John Dryden}}
\section{Introduction}
A sentence.

Another one.
Fascinating.

\subsection{One subheading}
Powered by \href{https://latex-ninja.com}{The \LaTeX{} Ninja blog}. What an elegant link. 

Of course, you can write in \textbf{bold} or ``quotes''.
Tilde (\textasciitilde) gives you a protected whitespace. \LaTeX{} even makes the effort to protect your personal space. Isn't that great? No, really, the thing is that \LaTeX{} will otherwise probably `eat' your spaces. Unlike MS Word which increases spacing inside words, \LaTeX{} increases spacing \emph{between} words and thus, handles whitespace pretty independently. A fair warning to you, you might have to endure Word user's comments along the lines of `I think you have an extra space here'. Whenever this happens, tell them that no, it's just good typesetting and that there is a self-help group which will help them get over it at \href{https://latex-ninja.com}{The \LaTeX{} Ninja blog}.\footcite[cf.][33--44]{CiteKey} % CiteKey is just whatever name you give to your source literature. Also, this comment won't show up in the output document - isn't that great? 

\begin{quote}\footnotesize
Towels can be harmful \dotdot in large quantities.\footcite[cf.][]{hitchhikersguide}
\end{quote}

\verb|\emph{}|, like the name kind of indicates, is used for emphasis. As opposed to \textit{textit} you can `layer' emphasis here.
\emph{so you can have \emph{emphasis} in your already emphasised block}. Cool, huh?
\texttt{source code} looks nicer in a typewriter font. But mind you, don't paste in source code which might annoy \LaTeX{}. If you're not sure, be on the safe side and use either
\begin{verbatim}
    the verbatim environment
\end{verbatim}
or else, the \verb|verb| command for short snippets.
\textsc{Authors} feel flattered if you put their names in small caps.

\begin{itemize}
    \item Of course, you can also take notes $\to$ very practical
\end{itemize}

\faBook~ This useless book symbol is meant to indicate that a to do list follows here:
\begin{enumerate}
    \item in case you need more law and order
    \item \verb|enumerate| environments (environment = all between \verb|\begin{}| and \verb|\end{}|)
    \begin{enumerate}
        \item can be nested.
    \end{enumerate}
\end{enumerate}

Certain characters, like \& have `internal meanings' in \LaTeX{}, so they have to be `escaped' (for example as \verb|\&|).
Generally, you don't need to change font sizes here at all since the semantic markup (declaring something to be a heading by using \verb|\section{Hint: this is a heading.}|). If you really feel like you still need to, it is encouraged that you only make changes using relative font {\footnotesize sizes}.

They offer the whole scale from {\scriptsize very small} to {\large large} or {\huge even larger}.\footcites{CiteKey,hitchhikersguide}

\section{Tables \& images}
\begin{wrapfigure}{R}{0.25\textwidth}
    \centering
    \includegraphics[width=0.15\textwidth]{test.png}
    \caption{A picture}
    \label{fig:the_pic}
\end{wrapfigure}
Like seen in [\ref{fig:the_pic}], graphics are extremely useful \parencite{articleInBook}.

\bigskip % gives you some additional space

\begin{figure}[h]
\begin{tabular}{c|l r}
    \textbf{What a} & \textbf{wonderful} & \textbf{table} \\ \hline
    centered & left-justified & right-justified  \\
    1 &  2 & 3
\end{tabular}
\end{figure}

\vspace{1cm} % reserves 1cm of extra space, works with em as well
Check the \href{http://mirror.utexas.edu/ctan/info/visualFAQ/visualFAQ.pdf}{Visual FAQ} for more hands-on info.

\section{Citations}

You can use this `manual' bibliography:

\begin{thebibliography}{99}

\bibitem{Principe} Lawrence R. \textbf{Principe}, {\it The Secrets of Alchemy} (Chicago 2013).

\bibitem{Comment} You can also just add stupid comments. Which hereby is officially discouraged by the author.

\end{thebibliography}

\subsection*{Using biblatex} % asterisk* will stop the heading from being numbered. It will not appear in the TOC

Or do it the automatic way (kind of like Citavi or Zotero). If you decide to take that path, there are many domain specific citation styles available (from the archaeological citation style of the German Archaeological Institute to all the standards (Chicago, etc.) and more. Just load it in the preamble where we declared how to use \texttt{biblatex}.

\printbibliography[title=My sources]
\newpage causes a manual page break. {\large\bfseries No you go and write that paper.}

Sorry to disappoint, but that process might take you quite a bit longer than 10 minutes. In case I come up with some ingenious productivity hack to enable you to write your seminar paper in 10 minutes, I'll let you know in another blog post on \href{https://latex-ninja.com}{The \LaTeX{} Ninja}.
\end{document}
